\documentclass{article}
\input{../../../../lib.sty}

\title{\texttt{NoisePrivacyMap<L2Distance<RBig>, ZeroConcentratedDivergence> for ZExpFamily<2>}}
\author{Michael Shoemate}
\date{}

\begin{document}

\maketitle

\contrib
Proves soundness of the implementation of \rustdoc{measurements/noise/trait}{NoisePrivacyMap} for \texttt{ZExpFamily<2>} in \asOfCommit{mod.rs}{f5bb719}.

\section{Hoare Triple}
\subsection*{Precondition}
\subsubsection*{Compiler-Verified}
\texttt{NoisePrivacyMap} is parameterized as follows:
\begin{itemize}
    \item \texttt{MI}, the input metric, is of type \rustdoc{metrics/type}{L2Distance<RBig>}
    \item \texttt{MO}, the output measure, is of type \rustdoc{measures/struct}{ZeroConcentratedDivergence}
\end{itemize}

\subsubsection*{User-Verified}
None

\subsection*{Pseudocode}
\lstinputlisting[language=Python,firstline=2,escapechar=|]{./pseudocode/NoisePrivacyMap_for_ZExpFamily2.py}

\subsection*{Postcondition}
\begin{theorem}
    Given a distribution \texttt{self},
    returns \texttt{Err(e)} if \texttt{self} is not a valid distribution.
    Otherwise the output is \texttt{Ok(privacy\_map)}
    where \texttt{privacy\_map} observes the following:

    Define \texttt{function(x) = x + Z} where \texttt{Z} is a vector of iid samples from \texttt{self}.

    For every pair of elements $x, x'$ in \texttt{VectorDomain<AtomDomain<IBig>{}>},
    and for every pair (\texttt{d\_in}, \texttt{d\_out}),
    where \texttt{d\_in} has the associated type for \texttt{input\_metric} and \texttt{d\_out} has the associated type for \texttt{output\_measure},
    if $x, x'$ are \texttt{d\_in}-close under \texttt{input\_metric}, \texttt{privacy\_map(d\_in)} does not raise an exception,
    and $\texttt{privacy\_map(d\_in)} \leq \texttt{d\_out}$,
    then \texttt{function(x)}, \texttt{function(x')} are \texttt{d\_out}-close under \texttt{output\_measure}.
\end{theorem}

\begin{proof}
    Line \ref{line:neg-scale} rejects \texttt{self} if \texttt{self} does not represent a valid distribution,
    satisfying the error conditions of the postcondition.

    We now construct the privacy map.
    First consider the extreme values of the scale and sensitivity parameters.
    The sensitivity \texttt{d\_in}, a bound on distances, must not be negative, as checked on line \ref{line:neg-sens}.
    In the case where sensitivity is zero (line \ref{line:zero-sens}), the privacy loss is zero, regardless the choice of scale parameter (even zero).
    This is because the privacy loss when adjacent datasets are always identical is zero.
    Otherwise, in the case where the scale is zero, the privacy loss is infinite.
    To avoid a rational division overflow, line \ref{line:zero-scale} returns infinity.

    By line \ref{line:map}, both the sensitivity and scale are positive rationals.
    Write $\sigma$ for \texttt{self.scale}.

    If the metric is non-modular, Theorem~14 from \cite{CKS20} implies that
    for all $\alpha > 1$,
    \[
        D_\alpha(M(x), M(x')) \leq \alpha \cdot \frac{(\texttt{d\_in}/\sigma)^2}{2}.
    \]

    Suppose instead that the metric is modular.
    Let
    \[
        G = \mathbb{Z}_{m_1} \times \cdots \times \mathbb{Z}_{m_d}
    \]
    be the quotient group induced by the coordinatewise moduli, and let
    $\pi : \mathbb{Z}^d \to G$ be coordinatewise reduction modulo those moduli.
    The wrapped mechanism is
    \[
        M_m(x) = \pi(x + Y),
    \]
    where the coordinates of $Y$ are independent discrete-Gaussian random variables with
    common scale $\sigma$.

    Consider neighboring inputs $x \sim x'$ and let $r = x' - x \in G$.
    By the definition of modular $L_2$ distance, there exists a lift
    $t \in \mathbb{Z}^d$ such that $\pi(t) = r$ and
    \[
        \sum_{i=1}^d \frac{t_i^2}{\sigma^2} \le \left(\frac{\texttt{d\_in}}{\sigma}\right)^2.
    \]

    Choose any lift $a \in \mathbb{Z}^d$ of $x$.
    Then $a + t$ is a lift of $x'$.
    Let $P$ be the law of $a + Y$ on $\mathbb{Z}^d$, and $Q$ the law of $a + t + Y$.
    Their pushforwards under $\pi$ are exactly the laws of $M_m(x)$ and $M_m(x')$.
    By data processing for R\'enyi divergence under the deterministic map $\pi$,
    \[
        D_\alpha(M_m(x), M_m(x')) \le D_\alpha(P, Q).
    \]
    Since adding the same constant does not change R\'enyi divergence,
    \[
        D_\alpha(P, Q) = D_\alpha(Y, Y + t).
    \]

    Proposition~5 of \cite{CKS20} gives the scalar discrete-Gaussian R\'enyi bound,
    and Theorem~14 gives the multivariate independent-coordinate version.
    Therefore, for all $\alpha > 1$,
    \[
        D_\alpha(Y, Y + t)
        \le
        \alpha \sum_{i=1}^d \frac{t_i^2}{2\sigma^2}
        \le
        \alpha \cdot \frac{(\texttt{d\_in}/\sigma)^2}{2}.
    \]
    Combining the displays yields the same bound for the wrapped mechanism.

    Line \ref{line:map} implements this upper bound exactly as
    \texttt{d\_in\textasciicircum 2 / (2 * scale\textasciicircum 2)} using exact rational arithmetic
    and a conservative cast to float.
\end{proof}

\bibliographystyle{alpha}
\bibliography{mod}

\end{document}
